% Source: physical PDF pp. 406--411 (printed pp. 383--388). Coverage:
% complete Section 22.4, from its heading through immediately before
% Section 22.5; Eqs. (22.4.1)--(22.4.3), Table 22.1, citations
% [11]--[15], no footnotes. Source title verified against the scan.
\subsection{Anomaly-Free Gauge Theories}
\label{sec:22.4}

We have calculated the effect of anomalies on the conservation of a general
current $J_a^\mu$. Where this current is itself coupled to a gauge
field, gauge invariance requires that the anomaly be absent. We saw in the
\hyperref[sec:22.3]{previous section} that the anomaly is proportional to
the completely symmetric constant factor $D_{abc}$ defined
by Eq.~\eqref{eq:22.3.12}, so for gauge currents we must
have~\hyperref[ch22-ref-11]{[11]}
\begin{equation}
 D_{abc}\equiv
 \frac{1}{2}\operatorname{tr}
 \bigl[\{T_a,T_b\}T_c\bigr]=0,
 \tag{22.4.1}\label{eq:22.4.1}
\end{equation}
where $T_a$ is the representation of the gauge algebra on the set of
all left-handed fermion and antifermion fields, and `$\operatorname{tr}$'
again denotes a sum over these fermion and antifermion species. This
condition may be satisfied for any gauge group if the fermion fields
furnish a suitable reducible or irreducible representation of the group.
In addition, there are some gauge groups for which Eq.~\eqref{eq:22.4.1}
is satisfied for fermions in \emph{any} representation of the
group~\hyperref[ch22-ref-12]{[12]}. (The Batalin--Vilkovisky formalism is
used in \hyperref[sec:22.6]{Section 22.6} to give a purely algebraic proof
that for such gauge groups the anomaly is absent to all orders of
perturbation theory.)

The condition \eqref{eq:22.4.1} is obviously satisfied if the left-handed
fermion (and antifermion) fields furnish a representation of the gauge
algebra that is equivalent to its complex conjugate, in the sense that
\[
 (\ii T_a)^*=S(\ii T_a)S^{-1}
\]
or equivalently (since we always take $T_a$ Hermitian)
\begin{equation}
 T_a^{\mathsf T}=-S T_a S^{-1}.
 \tag{22.4.2}\label{eq:22.4.2}
\end{equation}
(Inserting Eq.~\eqref{eq:22.4.2} in Eq.~\eqref{eq:22.3.12} gives
$D_{abc}=-D_{abc}$.) Such a representation
may be either \emph{real}, in which case it is possible by a similarity
transformation $T'_a=R T_a R^{-1}$ to convert the representation
to a form in which $T'_a$ is imaginary and antisymmetric, or
\emph{pseudoreal}, in which case this is impossible. (For instance, the
three-dimensional irreducible representation of $SU(2)$ is real, while
the two-dimensional representation is pseudoreal.) There is no anomaly
for gauge algebras that have only real or pseudoreal representations,
namely~\hyperref[ch22-ref-13]{[13]} $SO(2n+1)$ (including
$SU(2)=SO(3)$), $SO(4n)$ for $n\geq2$, $USp(2n)$ for $n\geq3$,
$G_2$, $F_4$, $E_7$, and $E_8$, and all of their direct sums. A few other
algebras also have only representations for which
$D_{abc}$ vanishes, even though some representations are
neither real nor pseudoreal~\hyperref[ch22-ref-12]{[12]}. These are
$SO(4n+2)$ (except for $SO(2)=U(1)$ and $SO(6)=SU(4)$) and $E_6$, and
their direct sums with each other and the above algebras. Anomalies are
thus only possible for gauge algebras that include $SU(n)$ (for $n\geq3$)
or $U(1)$ factors. As it happens, these are among the most important gauge
algebras in today's physics. The standard model is based on the gauge
group $SU(3)\times SU(2)\times U(1)$, so we must rely on cancellations among
the quarks and leptons to make the theory free of anomalies.

Table~\ref{tab:22.1} gives a classification of the left-handed spinor fields
of the first generation of the standard model according to the
representations they furnish of the color $SU(3)$ group and the electroweak
$SU(2)$ group, and the value of the $U(1)$ quantum number
$y/g'=t_3/g-q/e$.

\begin{table}[htbp]
 \centering
 \begingroup
 \renewcommand{\thetable}{22.1}
 \caption{First-generation left-handed fermion and antifermion fields of
 the standard model.}
 \label{tab:22.1}
 \renewcommand{\arraystretch}{1.35}
 \begin{tabular}{c@{\qquad}c@{\qquad}c@{\qquad}c}
  \hline\hline
  Fermions & $SU(3)$ & $SU(2)$ & $U(1)\ [y/g']$\\
  \hline
  $\begin{pmatrix}u\\ \dd\end{pmatrix}_{\!L}$
    & $\mathbf 3$ & $\mathbf 2$ & $-1/6$\\
  $u_R^*$ & $\overline{\mathbf 3}$ & $\mathbf{1}$ & $+2/3$\\
  $d_R^*$ & $\overline{\mathbf 3}$ & $\mathbf{1}$ & $-1/3$\\
  $\begin{pmatrix}\nu_e\\ e\end{pmatrix}_{\!L}$
    & $\mathbf{1}$ & $\mathbf 2$ & $1/2$\\
  $e_R^*$ & $\mathbf{1}$ & $\mathbf{1}$ & $-1$\\
  \hline\hline
 \end{tabular}
 \endgroup
\end{table}

We can now check whether $D_{abc}$ vanishes when $T_a$,
$T_b$, and $T_c$ run over all of the generators of
$SU(3)\times SU(2)\times U(1)$. We need only consider those combinations of
generators for which the product of $T_a$, $T_b$, and $T_c$
is neutral under $SU(3)\times SU(2)\times U(1)$, since
$D_{abc}$ obviously vanishes for all the others. We can make
invariants out of either zero, two, or three $SU(3)$ generators (because
$\mathbf8\times\mathbf8$ and
$\mathbf8\times\mathbf8\times\mathbf8$ both contain singlets), zero, two, or
three $SU(2)$ generators, and any number of $U(1)$ generators, so we only
need to check the following cases:

\noindent\textbf{\boldmath$SU(3)$--$SU(3)$--$SU(3)$:}
Here $D_{abc}$ vanishes because the left-handed fermions
furnish a representation
$\mathbf3+\mathbf3+\overline{\mathbf3}+\overline{\mathbf3}
+\mathbf{1}+\mathbf{1}+\mathbf{1}$ of $SU(3)$, which is real.

\noindent\textbf{\boldmath$SU(3)$--$SU(3)$--$U(1)$:}
Here the anomaly is proportional to
\[
 \sum_{\mathbf3,\overline{\mathbf3}}\frac{y}{g'}
 =-\frac{1}{6}-\frac{1}{6}+\frac{2}{3}-\frac{1}{3}=0.
\]

\noindent\textbf{\boldmath$SU(2)$--$SU(2)$--$SU(2)$:}
There is no anomaly here because $SU(2)$ only has real or pseudoreal
representations.

\noindent\textbf{\boldmath$SU(2)$--$SU(2)$--$U(1)$:}
Here the anomaly is proportional to
\[
 \sum_{\text{doublets}}\frac{y}{g'}
 =3\left(-\frac{1}{6}\right)+\frac{1}{2}=0.
\]

\noindent\textbf{\boldmath$U(1)$--$U(1)$--$U(1)$:}
Here the anomaly is proportional to
\[
 \sum\left(\frac{y}{g'}\right)^3
 =6\left(-\frac{1}{6}\right)^3
 +3\left(\frac{2}{3}\right)^3
 +3\left(-\frac{1}{3}\right)^3
 +2\left(\frac{1}{2}\right)^3
 +(-1)^3=0.
\]
We see that all anomalies cancel for the gauge symmetries of the standard
model~\hyperref[ch22-ref-13a]{[13a]}. This result can be neatly understood
by noting that $SU(3)\times SU(2)\times U(1)$ may be embedded in
$SO(10)$~\hyperref[ch22-ref-14]{[14]}. All of the representations of
$SO(10)$ are anomaly-free, so the same property is inherited by any
reducible representation of $SU(3)\times SU(2)\times U(1)$ that furnishes a
complete representation of $SO(10)$. As it happens, the left-handed fields
of a single generation of quarks, leptons, antiquarks, and antileptons plus
one additional $(SU(3)\times SU(2)\times U(1))$-singlet forms a complete
16-dimensional representation of $SO(10)$ (the fundamental spinor
representation), so for this set of left-handed fermions there are no
$SU(3)\times SU(2)\times U(1)$ anomalies. The singlet would not contribute
to such anomalies anyway, so there are no anomalies in the gauge
symmetries of the standard model.

There is one more anomaly that needs to be evaluated. All species of
fermions interact with gravitation in the same way. Calculation of the
fermion loop graph for the expectation value of the current
$\overline\chi T_a\gamma^\mu\chi$ (where $\chi$ is a column of all
left-handed fermion and antifermion fields) in the presence of an external
gravitational field yields an anomaly~\hyperref[ch22-ref-15]{[15]}
$\partial_\mu\langle\overline\chi T_a\gamma^\mu\chi\rangle$
proportional to
\[
 \operatorname{tr}\{T_a\}\epsilon^{\mu\nu\rho\sigma}
 R_{\mu\nu\kappa\lambda}R_{\rho\sigma}{}^{\kappa\lambda}.
\]
In particular, to avoid a gravitational violation of a gauge symmetry like
\eqref{eq:22.3.3}, the generators must satisfy
\begin{equation}
 \operatorname{tr}\{T_a\}=0.
 \tag{22.4.3}\label{eq:22.4.3}
\end{equation}
Like the pure gauge anomaly, this vanishes for gauge generators satisfying
Eq.~\eqref{eq:22.4.2}, so this anomaly receives no contribution from
fermions that form real or pseudoreal representations of the gauge group,
and it can therefore be calculated taking into account only those fermions
whose masses arise from a breaking of the gauge symmetry. Also, this
condition is automatically satisfied by the generators of any
\emph{simple} subalgebras like $SU(2)$ or $SU(3)$ of the gauge algebra.
($\operatorname{tr}\{T_a\}$ is just a number which commutes with all
the $T_b$, so if it is non-zero then the algebra is not simple.) Hence
we only need to check that Eq.~\eqref{eq:22.4.3} is satisfied by the $U(1)$
generators of the gauge algebra. In the standard model, the sum of the
values of the weak hypercharge $y$ for all left-handed fermions is
\[
 \sum\frac{y}{g'}
 =6\left(-\frac{1}{6}\right)
 +3\left(\frac{2}{3}\right)
 +3\left(-\frac{1}{3}\right)
 +2\left(\frac{1}{2}\right)+(-1)=0,
\]
so there are no gravitational anomalies in the standard model currents.

The requirement of vanishing of anomalies may be used as a guide in
formulating realistic theories. For instance, the values of the weak
hypercharges $y$ for various $SU(3)\times SU(2)$ multiplets were originally
taken from experiment, but one may wonder why these weak hypercharges (and
the corresponding average electric charges in each multiplet) take the
observed values. To answer this question, suppose we assign arbitrary weak
hypercharges $a$, $b$, $c$, $d$, and $e$ to the multiplets
$(u_L,d_L)$, $u_R^*$, $d_R^*$, $(\nu_L,e_L)$, and $e_R^*$,
respectively. The conditions for anomaly cancellation tell us that:

\noindent\textbf{\boldmath$SU(3)$--$SU(3)$--$U(1)$:}
\[
 \sum_{\mathbf3,\overline{\mathbf3}}y=2a+b+c=0;
\]
\noindent\textbf{\boldmath$SU(2)$--$SU(2)$--$U(1)$:}
\[
 \sum_{\text{doublets}}y=3a+d=0;
\]
\noindent\textbf{\boldmath$U(1)$--$U(1)$--$U(1)$:}
\[
 \sum y^3=6a^3+3b^3+3c^3+2\dd^3+e^3=0;
\]
\noindent\textbf{graviton--graviton--\boldmath$U(1)$:}
\[
 \sum y=6a+3b+3c+2d+e=0.
\]
Aside from the possibility of interchanging $u_R$ and $d_R$, these
equations have only two solutions, which we may call $U(1)$ and $U(1)'$:
\[
 \begin{array}{c@{\qquad}l}
 U(1):&
 b/a=-4,\quad c/a=2,\quad d/a=-3,\quad e/a=6,\\[2pt]
 U(1)':&
 b=-c,\quad a=c=d=e=0.
 \end{array}
\]
Furthermore, these solutions are exclusive; we cannot suppose that
\emph{both} $U(1)$ and $U(1)'$ are local symmetries, because then we would
encounter a $U(1)'$--$U(1)'$--$U(1)$ anomaly proportional to
$(-4)+(+2)\neq0$ and a $U(1)'$--$U(1)$--$U(1)$ anomaly proportional to
$(-4)^2-(+2)^2\neq0$. The $U(1)$ generator is just the weak hypercharge of
the standard electroweak theory (with the overall constant factor $a$
absorbed into the definition of $g'$), while the $U(1)'$ symmetry resembles
nothing observed in nature. This little calculation provides a rational
explanation of the assignment of $y$ values, or equivalently of electric
charges, in the standard model, and it shows that if all gauge anomalies
must cancel within a single generation of quarks and leptons, then it is
not possible to couple a gauge boson to any other $U(1)$ quantum number,
in addition to weak hypercharge.

On the other hand, although it is reasonable to guess that the
$SU(3)\times SU(2)\times U(1)$ gauge bosons of the standard model may couple
only to the known quarks and leptons (both to explain why no other fermions
have been discovered and to preserve the beautiful cancellation of
anomalies in the standard model), there may be other $U(1)'$ gauge bosons
that couple to other undetected $(SU(3)\times SU(2)\times U(1))$-neutral
fermions as well as to the known quarks and leptons. Suppose we denote the
$U(1)'$ quantum numbers $y'$ of the multiplets $(u_L,d_L)$, $u_R^*$,
$d_R^*$, $(\nu_L,e_L)$, and $e_R^*$ as $a'$, $b'$, $c'$, $d'$, and $e'$,
respectively. Since we don't know anything about possible
$(SU(3)\times SU(2)\times U(1))$-neutral fermions, the requirement of
cancellation of the $U(1)'$--$U(1)'$--$U(1)'$ and
graviton--graviton--$U(1)'$ anomalies does not help us to constrain $a'$,
$b'$, $c'$, $d'$, or $e'$. The remaining conditions of anomaly
cancellation tell us that:

\noindent\textbf{\boldmath$SU(3)$--$SU(3)$--$U(1)'$:}
\[
 \sum_{\mathbf3,\overline{\mathbf3}}y'=2a'+b'+c'=0;
\]
\noindent\textbf{\boldmath$SU(2)$--$SU(2)$--$U(1)'$:}
\[
 \sum_{\text{doublets}}y'=3a'+d'=0;
\]
\noindent\textbf{\boldmath$U(1)$--$U(1)$--$U(1)'$:}
\[
 \sum y^2y'
 =6a^2a'+3(-4)^2b'+3(2)^2c'+2(-3)^2d'+(6)^2e'=0;
\]
\noindent\textbf{\boldmath$U(1)$--$U(1)'$--$U(1)'$:}
\[
 \sum y{y'}^2
 =6{a'}^2+3(-4){b'}^2+3(2){c'}^2
 +2(-3){d'}^2+(6){e'}^2=0.
\]
The general solution has $y'$ a linear combination of $y$ and a quantum
number $B-L$ (where $B$ and $L$ are the conventional baryon and lepton
numbers) that takes values $1/3$, $-1/3$, $-1/3$, $-1$, and $+1$ for the
multiplets $(u_L,d_L)$, $u_R^*$, $d_R^*$, $(\nu_L,e_L)$, and $e_R^*$,
respectively. If $B-L$ is a local symmetry, with a coupling that is not
many orders of magnitude smaller than $e$, then it must be spontaneously
broken, since ordinary bodies have macroscopic values of $B-L$. To avoid
conflict with observations of neutral currents, the characteristic scale
$F$ of the symmetry breaking would have to be larger than that of the
electroweak interactions, but not necessarily many orders of magnitude
larger. Thus a neutral vector boson somewhat heavier than the $Z^0$ and
coupled to $B-L$ seems like the most plausible addition to the standard
model.

This has all been for a single generation of the standard model. With
three generations there are many more anomaly-free symmetries. One class
of symmetries that are not broken by anomalies or (as far as we know) by
anything else consists of the differences in the numbers of leptons of the
various flavors. Together with $B-L$, this will be important in
\hyperref[sec:23.5]{Section 23.5} in classifying the baryon- and
lepton-non-conserving processes produced by anomalies.
